\documentclass[11pt,a4paper]{article}

% Packages
\usepackage[utf8]{inputenc}
\usepackage[T1]{fontenc}
\usepackage{lmodern}
\usepackage[margin=2.5cm]{geometry}
\usepackage{graphicx}
\usepackage{booktabs}
\usepackage{longtable}
\usepackage{array}
\usepackage{xcolor}
\usepackage{hyperref}
\usepackage{caption}
\usepackage{float}
\usepackage{amsmath}
\usepackage{enumitem}
\usepackage{tabularx}
\usepackage{makecell}
\usepackage{titlesec}
\usepackage{fancyhdr}
\usepackage{pdflscape}

% Colors
\definecolor{darkblue}{HTML}{1B3A5C}
\definecolor{accentred}{HTML}{C0392B}
\definecolor{lightgray}{HTML}{F2F2F2}

% Hyperref setup
\hypersetup{
    colorlinks=true,
    linkcolor=darkblue,
    urlcolor=darkblue,
    citecolor=darkblue
}

% Header/Footer
\pagestyle{fancy}
\fancyhf{}
\fancyhead[L]{\small\textcolor{darkblue}{Indonesia Data Center Market Analysis}}
\fancyhead[R]{\small\textcolor{darkblue}{March 2026}}
\fancyfoot[C]{\thepage}
\renewcommand{\headrulewidth}{0.4pt}

% Title formatting
\titleformat{\section}{\Large\bfseries\color{darkblue}}{}{0em}{}[\titlerule]
\titleformat{\subsection}{\large\bfseries\color{darkblue}}{}{0em}{}

\begin{document}

% Title Page
\begin{titlepage}
\centering
\vspace*{3cm}
{\Huge\bfseries\textcolor{darkblue}{Data Centers in Indonesia}\\[0.5cm]}
{\LARGE\textcolor{accentred}{Market Landscape, Investment Database,\\and Fiscal Cost of Tax Incentives}\\[2cm]}
{\Large March 2026\\[1cm]}
\vfill
{\large Working Paper}\\[0.5cm]
{\small Compiled from public sources including Arizton, Mordor Intelligence,\\
company press releases, and Indonesian Ministry of Finance data.}
\end{titlepage}

\tableofcontents
\newpage

% ============================================================
\section{Executive Summary}
% ============================================================

Indonesia's data center market is experiencing rapid expansion, driven by the country's growing digital economy, hyperscaler cloud region deployments, and AI infrastructure demand. Key findings:

\begin{itemize}[leftmargin=*]
    \item \textbf{88 existing facilities} across 17 cities, with \textbf{25 upcoming} facilities in various stages of development.
    \item Total operational IT load capacity of approximately \textbf{307~MW} (2024), projected to reach \textbf{3,560~MW by 2030}.
    \item Over \textbf{USD~11.2~billion} in announced data center investments, led by Digital Edge (USD~4.5B), DAMAC Digital (USD~2.3B), and Microsoft (USD~1.7B).
    \item Indonesia offers \textbf{100\% corporate income tax holidays} for up to 20 years for qualifying data center investments above IDR~500~billion.
    \item We estimate the annual fiscal cost (revenue foregone) from data-center-specific tax holidays at approximately \textbf{USD~370~million per year}, with a cumulative ten-year cost of \textbf{USD~3.7~billion}.
\end{itemize}

% ============================================================
\section{Market Overview}
% ============================================================

\subsection{Market Size and Growth}

The Indonesia data center market was valued at \textbf{USD~2.81~billion in 2025} and is projected to reach \textbf{USD~6.08~billion by 2031}, growing at a CAGR of 13.73\%. In IT load capacity terms, the market is expected to expand from 1,440~MW in 2025 to 3,560~MW by 2030 (CAGR: 19.89\%).

\begin{table}[H]
\centering
\caption{Indonesia Data Center Market Projections}
\begin{tabular}{lrr}
\toprule
\textbf{Metric} & \textbf{2025} & \textbf{2030/2031} \\
\midrule
Market Value (USD billion)       & 2.81  & 6.08 (2031) \\
IT Load Capacity (MW)            & 1,440 & 3,560 (2030) \\
Construction Market (USD billion)& 3.05  & 7.11 (2030) \\
Colocation Segment (USD billion) & 0.46  & 1.15 (2030) \\
Hyperscale Segment (USD billion) & 3.49  & 7.96 (2030) \\
\bottomrule
\end{tabular}
\end{table}

\subsection{Geographic Distribution}

Jakarta dominates with \textbf{56.72\%} of market share in 2025. Batam is the fastest-growing secondary hub (CAGR: 21.70\%), benefiting from Special Economic Zone incentives, proximity to Singapore, and new submarine cable links.

\begin{table}[H]
\centering
\caption{Key Data Center Hubs in Indonesia}
\begin{tabular}{llr}
\toprule
\textbf{Region} & \textbf{Key Locations} & \textbf{Estimated Pipeline (MW)} \\
\midrule
Greater Jakarta & Bekasi, Cikarang, Cibitung, Bogor & $\sim$450 \\
Batam           & Nongsa Digital Park, SEZ           & $\sim$170 \\
Other           & Surabaya, Bandung, Pekanbaru       & $\sim$50 \\
\bottomrule
\end{tabular}
\end{table}

% ============================================================
\section{Data Center Inventory}
% ============================================================

\subsection{Major Operators by Capacity}

\begin{longtable}{p{3.8cm}p{3.5cm}rrp{3cm}}
\caption{Major Data Center Operators in Indonesia (Known Capacity)} \\
\toprule
\textbf{Operator} & \textbf{Key Facility} & \textbf{Live MW} & \textbf{Planned MW} & \textbf{Notes} \\
\midrule
\endfirsthead
\toprule
\textbf{Operator} & \textbf{Key Facility} & \textbf{Live MW} & \textbf{Planned MW} & \textbf{Notes} \\
\midrule
\endhead
\bottomrule
\endfoot

DCI Indonesia         & Cibitung Campus    & 119  & 600  & Largest operator; Tier~IV \\
Princeton Digital     & JC3 Jakarta        & 120  & 118  & AI-ready; Batam expansion \\
EdgeConneX            & Jakarta Campus     & 7    & 200  & Hyperscale multi-phase \\
Digital Edge          & EDGE1/2 + CGK      & 29   & 500  & \$4.5B CGK Campus \\
NTT DATA              & JKT2/JKT3/JKT2A   & 24.6 & 57   & JKT3 expandable to 45~MW \\
Telkom (NeutraDC)     & 33 sites           & 42   & 60   & 70\% utilization \\
SpaceDC               & ID01               & 25.5 & ---  & PUE 1.3 \\
STT GDC               & STT Jakarta        & ---  & 72   & Ground broken \\
BDx Indonesia         & CGK4 Jatiluhur     & ---  & 500  & AI campus; renewable \\
DAMAC Digital         & Cikarang           & ---  & 144  & \$2.3B; AI-ready \\
Digital Realty        & CGK10/CGK11        & 5    & 32   & JV with Bersama Digital \\
Gaw / Sinar Primera  & Batam              & 5.2  & 20   & Phase 1 operational \\
DayOne / INA          & Nongsa Digital Park& ---  & 72   & JV with INA \\
Equinix / Astra      & JK1 Jakarta        & ---  & ---  & AI-ready; 50+ NSPs \\
IndoKeppel            & IKDC1 Bogor        & ---  & ---  & 7-hectare campus \\
\end{longtable}

\subsection{Hyperscaler and Cloud Provider Presence}

\begin{table}[H]
\centering
\caption{Hyperscaler Investments in Indonesia}
\begin{tabular}{p{3cm}p{4cm}rp{4cm}}
\toprule
\textbf{Company} & \textbf{Project} & \textbf{Investment} & \textbf{Status} \\
\midrule
Microsoft     & Indonesia Central Region   & USD 1.7B & Under construction (2024--2028) \\
AWS           & Multi-AZ Jakarta Region    & ---      & Operational \\
Google Cloud  & Jakarta Cloud Region       & ---      & Operational \\
Indosat-NVIDIA& AI Factory Jakarta         & USD 250M & Operational (Oct 2024) \\
Tencent       & Capacity Expansion         & USD 500M & Planned (2026) \\
Oracle        & Indonesia North (Batam)    & ---      & Operational (Jul 2025) \\
\bottomrule
\end{tabular}
\end{table}

% ============================================================
\section{Investment Landscape}
% ============================================================

\subsection{Announced Investments Summary}

Total announced data center investments in Indonesia exceed \textbf{USD~11.2~billion}. The following table summarizes the largest commitments:

\begin{table}[H]
\centering
\caption{Top Data Center Investments by Value}
\begin{tabular}{p{3.5cm}rrp{4cm}}
\toprule
\textbf{Company} & \textbf{USD (B)} & \textbf{MW} & \textbf{Timeline} \\
\midrule
Digital Edge       & 4.50 & 500  & 2026--2027 (Phase 1) \\
DAMAC Digital      & 2.30 & 144  & 2026 \\
Microsoft          & 1.70 & ---  & 2024--2028 \\
JV Consortia       & 0.75 & ---  & 2024--2026 \\
Tencent            & 0.50 & ---  & 2026 \\
Digital Edge (EDGE2)& 0.33 & 23  & 2026 \\
Indosat-NVIDIA     & 0.25 & 80   & 2024 \\
Princeton Digital  & 0.11 & 118  & 2025--2026 \\
\midrule
\textbf{Total}     & \textbf{$\sim$10.4} & & \\
\bottomrule
\end{tabular}
\end{table}

\noindent Additional unlisted investments (BDx 500~MW campus, DCI Indonesia expansions, EdgeConneX 200~MW, STT GDC, and others) bring the total well above USD~11~billion.

% ============================================================
\section{Tax Incentive Framework}
% ============================================================

\subsection{Corporate Income Tax Holiday}

Under \textbf{PMK No.~69/2024}, data centers are classified as one of 18 \textit{pioneer industries} eligible for corporate income tax (CIT) holidays:

\begin{table}[H]
\centering
\caption{Tax Holiday Structure for Data Center Investments}
\begin{tabular}{p{5cm}rr}
\toprule
\textbf{Investment Size} & \textbf{CIT Reduction} & \textbf{Duration} \\
\midrule
IDR 100B--500B (USD 6.6--33.3M)  & 50\%  & 5 years \\
IDR 500B--1T                      & 100\% & 5 years \\
IDR 1T--5T                        & 100\% & 7 years \\
IDR 5T--15T                       & 100\% & 10 years \\
IDR 15T--30T                      & 100\% & 15 years \\
Above IDR 30T (USD $\sim$2B)      & 100\% & 20 years \\
\midrule
\multicolumn{3}{l}{\textit{+ 50\% CIT reduction for 2 additional years after holiday ends}} \\
\bottomrule
\end{tabular}
\end{table}

\subsection{Special Economic Zone (SEZ) Incentives}

Facilities in designated SEZs (Batam, Cikarang, Nongsa Digital Park) benefit from:
\begin{itemize}[leftmargin=*]
    \item \textbf{0\% VAT} on imported equipment (saving $\sim$11\% capex)
    \item \textbf{100\% foreign ownership} permitted
    \item \textbf{Accelerated depreciation} on digital infrastructure assets
    \item Streamlined \textbf{10-week approval} via Online Single Submission (OSS), down from 24 weeks
\end{itemize}

\subsection{Other Fiscal Incentives}

\begin{itemize}[leftmargin=*]
    \item \textbf{Investment allowance}: 30\% of total investment deductible over 6 years (5\%/year)
    \item \textbf{Reduced dividend WHT}: 10\% (vs.\ standard 20\%)
    \item \textbf{Tax loss carry-forward}: up to 10 years (vs.\ standard 5)
    \item \textbf{R\&D super deduction}: up to 300\% of qualifying expenses
\end{itemize}

\subsection{OECD Pillar Two Implications}

The 15\% global minimum tax creates a structural challenge: if Indonesia grants a full CIT holiday, the investor's home country may levy a top-up tax, effectively transferring the fiscal benefit abroad. Indonesia is redesigning incentives to incorporate a \textbf{domestic top-up tax} (Qualified Domestic Minimum Top-up Tax, or QDMTT), ensuring the 15\% floor is collected domestically rather than ceded to foreign jurisdictions.

% ============================================================
\section{Fiscal Cost Assessment}
% ============================================================

\subsection{Indonesia's Overall Tax Expenditure}

Indonesia's total tax expenditure in 2023 was \textbf{IDR~362.5~trillion (USD~23.8~billion)}, equivalent to 1.73\% of GDP. Tax holidays and allowances specifically accounted for IDR~20~trillion over the 2018--2022 period, attracting IDR~370~trillion in investment---an \textbf{18.5x leverage ratio}.

\subsection{Estimated Data Center Tax Revenue Foregone}

We construct a bottom-up estimate of the fiscal cost of tax holidays for the data center sector:

\begin{table}[H]
\centering
\caption{Estimated Annual Revenue Foregone from Data Center Tax Holidays}
\begin{tabular}{lr}
\toprule
\textbf{Parameter} & \textbf{Value} \\
\midrule
Total announced DC investment        & USD 11.2 billion \\
Assumed annual revenue (at maturity)  & USD 7.6 billion \\
Assumed operating profit margin       & 22\% \\
Implied taxable income                & USD 1.68 billion \\
CIT rate                              & 22\% \\
\textbf{Annual CIT foregone (100\% holiday)}  & \textbf{USD 370 million} \\
\textbf{Equivalent in IDR}                    & \textbf{$\sim$IDR 5.7 trillion} \\
\midrule
\textit{10-year cumulative foregone}  & \textit{USD 3.7 billion} \\
\textit{As \% of 2023 total tax expenditure}  & \textit{$\sim$1.6\%} \\
\bottomrule
\end{tabular}
\end{table}

\noindent\textbf{Key assumptions and caveats}:
\begin{enumerate}[leftmargin=*]
    \item Not all investments will qualify for the maximum 100\% holiday; some receive 50\% reductions.
    \item Data centers have long ramp-up periods (2--4 years) before reaching full profitability; revenue foregone is lower in early years.
    \item OECD Pillar Two will reduce the effective holiday for large MNEs by imposing a 15\% domestic top-up, recovering approximately \textbf{USD~250~million annually} of the foregone amount.
    \item The estimate does not include \textbf{indirect fiscal benefits}: job creation, digital economy growth (target: USD~130B by 2025), increased VAT from downstream services, and property/land taxes.
    \item SEZ-related VAT exemptions on equipment imports add an estimated \textbf{USD~120--180~million} in additional revenue foregone (one-time, during construction phases).
\end{enumerate}

\subsection{Cost-Benefit Perspective}

\begin{table}[H]
\centering
\caption{Fiscal Cost vs.\ Economic Benefits (Indicative)}
\begin{tabular}{p{6cm}r}
\toprule
\textbf{Metric} & \textbf{Estimate} \\
\midrule
Annual tax revenue foregone (CIT)             & USD 370M \\
SEZ VAT foregone (construction, one-time)     & USD 120--180M \\
Direct jobs created (construction + ops)       & 15,000--25,000 \\
Indirect/induced employment                   & 50,000--80,000 \\
Data center market value contribution (2025)   & USD 2.81B \\
Digital economy GDP contribution target        & USD 130B \\
Investment leverage ratio (historical)         & 18.5x \\
\bottomrule
\end{tabular}
\end{table}

% ============================================================
\section{Conclusions}
% ============================================================

\begin{enumerate}[leftmargin=*]
    \item Indonesia has become a \textbf{major data center destination} in Southeast Asia, with over USD~11B in announced investments and 25 facilities under development.
    \item The tax holiday regime is a \textbf{significant driver} of investment: historical data shows an 18.5x investment-to-incentive leverage ratio.
    \item The estimated annual fiscal cost of \textbf{USD~370M} (IDR~5.7T) is substantial but modest relative to Indonesia's total tax expenditure (IDR~362.5T) and the economic multiplier effects.
    \item The \textbf{OECD Pillar Two} global minimum tax will structurally reduce the effective cost of tax holidays by ensuring a 15\% floor, potentially recovering USD~250M annually.
    \item \textbf{Jakarta} will remain the primary hub, but \textbf{Batam}'s SEZ advantages position it as a fast-growing alternative, especially for Singapore-linked workloads.
\end{enumerate}

% ============================================================
\section*{Data Sources}
% ============================================================

\begin{itemize}[leftmargin=*]
    \item Arizton Advisory \& Intelligence --- Indonesia Data Center Portfolio (2025--2028)
    \item Mordor Intelligence --- Indonesia Data Center Market Report (2025--2031)
    \item ResearchAndMarkets.com --- Indonesia Colocation Portfolio Analysis
    \item PWC Tax Summaries --- Indonesia Corporate Tax Credits and Incentives
    \item ASEAN Briefing --- Tax Incentives for Businesses in Indonesia
    \item Jakarta Globe --- Indonesia Extends Tax Holiday to 2026
    \item Indonesian Ministry of Finance (BKF) --- Tax Expenditure Report 2023
    \item Company press releases: Digital Edge, DAMAC Digital, NTT DATA, Princeton Digital Group, DCI Indonesia, Microsoft, EdgeConneX
    \item DataCenterKnowledge.com, DataCenterDynamics.com, DataCenterFrontier.com
\end{itemize}

\end{document}
